% ------------------------------------------------------------
% §6  Differential cryptanalysis: the adversary's calculus
% ------------------------------------------------------------
\section{The adversary's calculus: differential cryptanalysis}
\label{sec:differential}

\Cref{sec:below} reported the adversarial state of the art as bare round
counts. This section opens the instrument that produces those numbers.
Differential cryptanalysis is the only mathematical technology that has
ever killed a member of this design family---it felled MD5, SHA-0, and
SHA-1 (\cref{sec:falling})---and its core idea is elementary enough to
state in one line: \emph{since the function cannot be understood, study
its finite differences instead}. What follows is the calculus itself,
its startling secret history, the story of its mechanization, and an
honest account of where, on \SHA{}, it is currently stuck. Readers who
absorb this section will understand precisely what the ``round counts''
of \cref{sec:below} measure.

\subsection{A calculus of differences}\label{sec:diffcalc}

For a map $F$ on bit strings and a fixed \emph{input difference}
$\Delta$, consider the discrete derivative
\[
  D_{\Delta}F(x) \;=\; F(x \shaxor \Delta) \,\shaxor\, F(x),
\]
the exact finite-difference analogue of differentiation, with $\shaxor$
playing subtraction in the $\Ftwo$-geometry of \cref{sec:statespace}.
The adversary's question is distributional: for random $x$, what does
$D_{\Delta}F(x)$ look like?

Two observations power everything. First, \emph{differencing kills
constants and linear structure}. If $L$ is $\Ftwo$-linear---a rotation,
a shift, one of the $\Sigma/\sigma$ maps of \cref{sec:roundshape}---then
\[
  D_{\Delta}L(x) \;=\; L(\Delta) \qquad \text{for every } x:
\]
the derivative is a \emph{constant}, independent of the unknown input,
exactly as $\mathrm{d}/\mathrm{d}x$ annihilates constants and
linearizes. All the values the attacker does not know---the other
message words, the chaining state, the round constants $K_t$---drop out
of the linear part of every round. This is why one differentiates: a
difference propagates through the transparent parts of the circuit
deterministically, leaving a residue of uncertainty only at the
nonlinear gates.

Second, at those nonlinear gates---the modular additions $\shaplus$ and
the quadratic functions $\mathrm{Ch}, \mathrm{Maj}$---a difference does
not propagate deterministically but \emph{branches into a probability
distribution over output differences}. The fundamental case is a single
addition (\cref{fig:diffadd}). Give one addend a one-bit difference at
position $i$ and add the same unknown $y$ to both members of the pair:
the two sums differ at bit $i$ certainly, but whether the difference
stops there is decided by the \emph{carries}, which the $\Ftwo$-lens
cannot see. If the carries out of position $i$ agree---this happens for
exactly half of all $(x,y)$---the output difference is again the single
bit $e_i$; otherwise a carry-difference runs leftward, one further bit
at a time, each extra bit halving the probability. One position is
exceptional: a difference in the top bit ($i = 31$) propagates with
probability $1$, because the disagreeing carry falls off the end of the
word and is discarded by the reduction mod $2^{32}$. Attackers
therefore park differences in the most significant bit whenever they
can---a free move in an economy where everything else costs. The
complete theory of differential propagation through addition (arbitrary
difference patterns, exact probabilities, all computable in linear
time) was worked out by \citet{LipmaaMoriai2001}, and its headline is a
pleasing cost principle: \emph{the probability of the straight-through
transition is $2^{-(\text{number of active bits, msb excluded})}$}. To
first order, difference weight is money, and the exchange rate is a
factor of two per active bit. The bitwise functions obey the same
economy: at each bit where its inputs differ, $\mathrm{Ch}$ or
$\mathrm{Maj}$ passes or absorbs the difference depending on the values
of the other, unknown operands---a coin flip per active bit, one more
factor of $\tfrac12$ each.

\begin{figure}[ht]
  \centering
  \begin{tikzpicture}[font=\small,
      cell/.style={draw, minimum width=5.2mm, minimum height=5.2mm,
                   inner sep=0pt},
      dbit/.style={cell, fill=red!25},
      qbit/.style={cell, fill=lenfill},
      lbl/.style={font=\scriptsize}]
    % ---------------- left panel: difference at interior bit 3 ------
    \foreach \i [evaluate=\i as \x using {(7-\i)*0.52}] in {0,...,7} {
      \ifnum\i=3
        \node[dbit] (a\i) at (\x,0) {$1$};
      \else
        \node[cell] (a\i) at (\x,0) {};
      \fi
      \node[lbl, above=0.5mm of a\i] {\i};
    }
    \node[lbl, anchor=east] at (-0.55,0) {$\Delta x$};
    \draw[->] (1.82,-0.45) -- (1.82,-1.25)
      node[midway, right=1mm, lbl, align=left]
      {$\shaplus\, y$ (same $y$ on both\\ members of the pair)};
    \foreach \i [evaluate=\i as \x using {(7-\i)*0.52}] in {0,...,7} {
      \ifnum\i=3
        \node[dbit] (b\i) at (\x,-1.7) {$1$};
      \else
        \ifnum\i>3 \ifnum\i<7
          \node[qbit] (b\i) at (\x,-1.7) {?};
        \else
          \node[cell] (b\i) at (\x,-1.7) {};
        \fi
        \else
          \node[cell] (b\i) at (\x,-1.7) {};
        \fi
      \fi
    }
    \node[lbl, anchor=east] at (-0.55,-1.7) {$\Delta(x \shaplus y)$};
    \draw[->, thick, orange!80!black]
      (b3.north west) ++(-0.05,0.06) to[bend right=20] (-0.15,-1.05);
    \node[lbl, orange!80!black, align=right, anchor=east] at (-0.55,-0.85)
      {carry difference\\ runs leftward};
    \node[lbl, align=left, anchor=north west] at (-0.6,-2.45)
      {stops immediately: probability $\tfrac12$\\
       each further bit: another $\tfrac12$};
    % ---------------- right panel: msb difference --------------------
    \begin{scope}[xshift=6.4cm]
    \foreach \i [evaluate=\i as \x using {(7-\i)*0.52}] in {0,...,7} {
      \ifnum\i=7
        \node[dbit] (c\i) at (\x,0) {$1$};
      \else
        \node[cell] (c\i) at (\x,0) {};
      \fi
      \node[lbl, above=0.5mm of c\i] {\i};
    }
    \node[lbl, anchor=east] at (-0.55,0) {$\Delta x$};
    \draw[->] (1.82,-0.45) -- (1.82,-1.25)
      node[midway, right=1mm, lbl] {$\shaplus\, y$};
    \foreach \i [evaluate=\i as \x using {(7-\i)*0.52}] in {0,...,7} {
      \ifnum\i=7
        \node[dbit] (d\i) at (\x,-1.7) {$1$};
      \else
        \node[cell] (d\i) at (\x,-1.7) {};
      \fi
    }
    \node[lbl, anchor=east] at (-0.55,-1.7) {$\Delta(x \shaplus y)$};
    \draw[->, thick, gray]
      (d7.north west) ++(-0.02,0.08) to[bend right=20] ++(-0.5,0.35);
    \node[lbl, gray, align=left] at (0.6,-0.75) {carry falls\\ off the word};
    \node[lbl, align=left, anchor=north west] at (-0.6,-2.45)
      {top bit: probability $1$\\ (the attacker's free move)};
    \end{scope}
  \end{tikzpicture}
  \caption{How a one-bit difference passes through modular addition
  (bit $7$ = most significant, in this $8$-bit cartoon of a $32$-bit
  word). \emph{Left:} an interior difference survives unchanged only if
  the carries agree---probability $\tfrac12$---and every extra bit of
  carry-difference costs another factor $\tfrac12$---the Lipmaa--Moriai
  cost formula of the text. \emph{Right:} a difference in the top bit
  always passes intact, since the disagreeing carry is discarded mod
  $2^{32}$. Differential cryptanalysis views the whole of \SHA{}
  through this lens: linear pieces are free wires; every active bit at
  a nonlinear gate levies a toll of $\tfrac12$.}
  \label{fig:diffadd}
\end{figure}

Notice what has happened in the language of \cref{sec:statespace}: the
attacker has \emph{chosen the $\Ftwo$-geometry} and agreed to pay, in
probability, at every point where the design invokes the other one. All
of the analytical cost is concentrated exactly at the carry chains and
quadratic gates---the terms $\Ftwo$-linear algebra cannot see. There is
a dual calculus using arithmetic differences $x' - x \bmod 2^{32}$,
which tames the additions and instead pays at the xors and rotations;
the craft of the great hash attacks (\cref{sec:craft}) was, in large
part, the discovery that one should track \emph{both} lenses
simultaneously, bit by bit, in the form of signed differences. The
two-geometry tension that \cref{sec:statespace} presented as the
design's defense is thus also, precisely, the structure of the attack:
differential cryptanalysis is what the two-lens problem looks like when
an adversary refuses to give up.

\subsection{Trails, tolls, and free rounds}\label{sec:craft}

A single round is only the first step. A \emph{differential
characteristic} (or \emph{trail}) prescribes a difference for every
intermediate quantity through all the rounds---a full flight plan for
the perturbation---and, under a heuristic independence assumption, its
probability is the \emph{product} of the per-gate transition
probabilities along the way. An attacker who wants a collision needs a
trail beginning at a nonzero message difference and ending, after the
feedforward, at state difference \emph{zero}; if the trail has
probability $p$, then about $1/p$ random conforming attempts are
expected before one message pair follows the entire flight plan and the
hashes agree. The attack economy is thus brutally simple: \emph{total
toll} $=$ sum, over all rounds, of active bits at nonlinear gates;
\emph{cost} $\approx 2^{\text{toll}}$; the design's job is to force
every possible flight plan to accumulate an unpayable toll.

\begin{figure}[ht]
  \centering
  \begin{tikzpicture}[font=\small,
      inj/.style={draw, fill=red!25, minimum width=8mm, minimum height=5mm,
                  font=\scriptsize},
      cor/.style={draw, fill=msgfill, minimum width=8mm, minimum height=5mm,
                  font=\scriptsize},
      non/.style={draw, fill=white, minimum width=8mm, minimum height=5mm},
      lbl/.style={font=\scriptsize}]
    % round axis
    \foreach \r/\x in {t/0, t+1/1.15, t+2/2.3, t+3/3.45, t+4/4.6, t+5/5.75}
      \node[lbl] at (\x, 2.15) {$\r$};
    \node[lbl, anchor=east] at (-0.9, 1.55) {message diff.\ $\delta W$};
    \node[inj] at (0,1.55) {seed};
    \foreach \x in {1.15, 2.3, 3.45, 4.6, 5.75}
      \node[cor] at (\x,1.55) {corr.};
    % state difference weight bars
    \node[lbl, anchor=east] at (-0.9, 0.6) {state diff.\ weight};
    \draw[->] (-0.55,0.05) -- (7.0,0.05);
    \fill[red!50]    (-0.2,0.05) rectangle (0.25,0.65);
    \fill[red!35]    (0.95,0.05) rectangle (1.4,0.5);
    \fill[red!35]    (2.1,0.05) rectangle (2.55,0.5);
    \fill[red!35]    (3.25,0.05) rectangle (3.7,0.42);
    \fill[red!25]    (4.4,0.05) rectangle (4.85,0.3);
    \draw[very thick, projframe] (5.55,0.05) -- (7.0,0.05)
      node[pos=0.6, above, lbl, projframe] {zero again};
    % toll ledger
    \node[lbl, anchor=east] at (-0.9,-0.55) {toll this round};
    \foreach \x/\p in {0/2^{-1}, 1.15/2^{-3}, 2.3/2^{-3}, 3.45/2^{-2},
                       4.6/2^{-2}, 5.75/2^{-1}}
      \node[lbl] at (\x,-0.55) {$\p$};
    \node[lbl, anchor=west] at (6.35,-0.55) {$\Rightarrow\ \prod = 2^{-12}$};
    \draw[decorate, decoration={brace, mirror, amplitude=5pt}]
      (-0.35,-0.95) -- (6.1,-0.95)
      node[midway, below=6pt, lbl]
      {one \emph{local collision}: inject, shepherd, cancel};
  \end{tikzpicture}
  \caption{Cartoon of a \emph{local collision}, the modular building
  block of every attack on this family since Chabaud and Joux's 1998
  analysis of SHA-0 (cited in the text): a seed difference is injected
  through one
  message word, then five successive message-word corrections shepherd
  and finally cancel the disturbance, leaving the state clean. Each
  round levies its toll (\cref{fig:diffadd}); the trail's probability
  is the product. A full attack chains many local collisions along a
  pattern compatible with the message schedule---and the toll of the
  best full trail is what decides whether the design lives or dies.}
  \label{fig:localcollision}
\end{figure}

Three refinements turn this arithmetic into the attacks of
\cref{sec:falling}, and all three are conceptually simple. First,
\emph{local collisions} (\cref{fig:localcollision}): rather than one
monolithic trail, the attacker builds a short pattern---inject a
difference through one message word, cancel it a few rounds later
through others---and then chains copies of the pattern wherever the
message schedule permits~\citep{ChabaudJoux1998}. Second, \emph{message
freedom}: for the first sixteen rounds the words $W_t$ are the message
itself (\cref{sec:roundshape}), so the attacker does not \emph{hope}
the early tolls are paid---she \emph{solves} for message bits that pay
them deterministically. Wang's celebrated breaks of MD5 and
SHA-1~\citep{WangYu2005, WangYinYu2005} rested on pushing this
``message modification'' astonishingly deep; the practical effect is
that the toll meter starts running only around round $17$, and the
first quarter of the function contributes nothing to its own defense.
Third, \emph{automation}: what Wang's school did by virtuosic hand
calculation---juggling xor- and arithmetic-differences as signed,
bit-level conditions---was recast by \citet{DeCanniereRechberger2006}
as a constraint-propagation search over ``generalized conditions,''
mechanically discovering trails no human would find. The current
records on \SHA{} are held by exactly this lineage of tools, now
built on MILP, \textsc{sat}, and \textsc{smt}
solvers~\citep{MendelNadSchlaffer2013, LiLiuWang2024}, closing the loop
with \cref{sec:npsearch}: the adversary's calculus has itself become a
branch of solver engineering.

\subsection{A method twenty years secret}\label{sec:secrethistory}

The published history begins in 1990, when Biham and Shamir announced
differential cryptanalysis and used it to attack reduced variants of
DES, the American encryption standard of
1977~\citep{BihamShamir1991}. A curiosity surfaced immediately: DES,
designed in the early 1970s, was \emph{strangely well prepared} for the
new attack. Its nonlinear tables were, component for component, close
to optimally resistant; the method that broke every academic variant
of DES bounced off the real thing.

The explanation arrived in 1994, and it is one of the most remarkable
documents in the literature. Coppersmith, of the original IBM design
team, was cleared to reveal the design criteria and wrote them down
plainly: the IBM team had discovered differential cryptanalysis
themselves---internally, the ``T-attack''---\emph{in 1974}, and
hardened DES against it before standardization; the NSA, consulted on
the design, asked that the rationale be kept secret, since publishing
the criteria ``would reveal the technique of differential
cryptanalysis, a powerful technique that can be used against many
ciphers''~\citep{Coppersmith1994}. The central tool of public
cryptanalysis was thus twenty years old on the day it was discovered.
For sixteen of those years, every open researcher who studied DES was
probing defenses calibrated against a weapon whose existence they did
not suspect.

This is not merely a good story; it is a load-bearing datum for the
epistemology of \cref{sec:epistemology}, because the pattern
\emph{repeats within our own object's lineage}. Recall the footnote of
\cref{sec:history}: in 1994--95 the NSA withdrew SHA-0 and changed a
single rotation, citing an undisclosed weakness; the public rediscovery
of what that tweak defends against---a differential collision
attack---took until 1998~\citep{ChabaudJoux1998}. Twice measured, the
gap between the classified and the public understanding of this exact
attack has been twenty years and four years. The honest inference for a
naturalist: the map of \cref{sec:below} is the \emph{public} frontier,
a lower bound on the true state of knowledge, drawn by a community that
has historically been behind by years to decades---and \SHA{}'s design
parameters were set by the party that was ahead. The comfort and the
discomfort of that fact are left to the reader to apportion.

\subsection{The siege of \SHA{}, and where it stands}\label{sec:siege}

Why did the calculus shatter MD5 and SHA-1 while \SHA{}, made of the
same primitives, holds? The anatomy of \cref{sec:roundshape} contains
the visible answers. SHA-1's message schedule is entirely
$\Ftwo$-linear, so schedule differences propagate by exact linear
algebra and \emph{sparse} schedules---low total weight, hence low
toll---can be found by searching a linear code for low-weight words;
this is precisely the handle the collision attacks
gripped (\cref{sec:hope}). \SHA{} interleaves modular additions into
the schedule itself and taxes every round twice over: each of
$\Sigma_0, \Sigma_1$ fans any single active bit into three, both tracks
of the two-track recurrence pass through a quadratic gate every round,
and the $\sigma$-mixing makes a sparse schedule difference bloom
within a few steps. The result, empirically: low-toll trails simply
stop existing a dozen-odd rounds past the free zone, and two decades of
search---human, then mechanical---have moved the frontier only from
about $19$ rounds (2006) to $31$ rounds (collisions) today. The precise
current records: practical colliding pairs for $28$ rounds and a
$2^{65.5}$ collision attack on $31$ rounds of the $64$, both from the
automated-search school~\citep{MendelNadSchlaffer2013}; practical
\emph{semi-free-start} collisions (the attacker may choose the chaining
value, removing the Merkle--Damg{\aa}rd anchor of \cref{thm:md}) for
$38$ rounds in 2013 and $39$ rounds in 2024, the latter found with a
\textsc{sat}/\textsc{smt} trail search~\citep{LiLiuWang2024,
AlamgirNejatiBright2024}; and, on the inversion side, preimages to $45$
rounds at a cost indistinguishable from brute
force~\citep{KRS2012}. Li, Liu, and Wang open their record-setting 2024
paper with the flat sentence that ``there is almost no progress in
collision attacks on SHA-2 after ASIACRYPT 2015''---a decade of solver
revolutions purchasing one round.

The shape of the stall is worth internalizing, because it is the
quantitative content of the phrase ``security margin.'' Rounds $1$--$16$
are free (message freedom); rounds $17$ to roughly $40$ are the
contested zone, where automated searches still find trails but each
round multiplies the toll; beyond that, \emph{no usable trail has ever
been exhibited}, not because a theorem forbids one but because every
known search method drowns. The design's $64$ rounds thus stand at
about twice the depth the best siege engines have ever reached with
unlimited choice of starting position. Whether that factor of two is a
fortress wall or a picket fence is exactly the question
\cref{sec:falling} showed nobody can answer from first principles: no
theory predicts the frontier's future velocity, and the SHA-1
precedent---from academic $2^{69}$ to industrial collision in twelve
years---is the reason nobody treats the margin as a moat. (One further
credential: the compression function's core, extracted and run as a
block cipher under the name SHACAL-2, survived the open European
NESSIE evaluation and was selected for its 2003 portfolio---the same
engine vetted from a second angle, as a cipher rather than a hash.)

For this paper's program the moral is double. The adversary's calculus
is the most refined body of exact knowledge about \emph{reduced}
members of the family---every trail is a theorem-with-certificate about
a specific truncation, machine-checkable pair in hand---and at the same
time it is purely \emph{local} knowledge: a trail certifies one orbit
of one difference, and says nothing global about the map. The
statistical coordinates of \cref{sec:randommaps} and the trail calculus
of this section are complementary instruments, and to our knowledge no
one has systematically pointed them at the same scaled-down family to
see whether the round where trails die is the round where the Flajolet
silhouette locks in (\cref{sec:projects}).

\begin{project}{Differential trails in miniature}{probability;
  programming; patience with bit manipulation}{4--6 weeks}
  Build the toll calculus for a toy \SHA{} on $8$-bit words. First
  verify \cref{fig:diffadd} empirically: tabulate the exact
  distribution of $\Delta(x \shaplus y)$ over all pairs for every
  one-bit $\Delta x$, and check the Lipmaa--Moriai
  cost formula~\citep{LipmaaMoriai2001}. Then implement a
  branch-and-bound search for the lowest-toll trail through $r$ rounds
  of the toy compression function and plot $\log_2(\text{best trail
  probability})$ against $r$. Two landmarks make the plot speak: the
  birthday line ($-n/2$ for an $n$-bit toy state), below which a trail
  is useless to an attacker, and the round at which your search itself
  starts to drown. You will have reproduced, at fruit-fly scale, both
  the designer's round-count decision and the attacker's wall---the two
  sides of \cref{sec:siege}---and the same code instruments the
  comparative questions of \cref{sec:projects}.
\end{project}

\subsection{Coda: can a machine learn the difference?}\label{sec:mlcrypt}

One newer thread deserves flagging, because it closes a loop with the
proof-barrier of \cref{sec:barriers} in a way that is either profound or
a curiosity, and nobody yet knows which. In 2019 Gohr trained a neural
network to distinguish the outputs of a round-reduced ARX cipher
(Speck, a near-relative of \SHA{}'s primitives) from random, and it
\emph{beat} the best hand-built differential distinguishers---and, on
inspection, had learned features \emph{beyond} the difference
distribution table, correlations the human theory of \cref{sec:diffcalc}
does not track~\citep{Gohr2019}. Read against Razborov--Rudich
(\cref{sec:barriers}), this is quietly startling. A trained
distinguisher \emph{is} a natural property in their precise sense---an
efficiently computable, large predicate that separates the function
from random---so a machine that reliably learns such properties on
\SHA{}-family maps would be walking, empirically, along the exact edge
the natural-proofs barrier says cannot be crossed at full scale.
Whether the network finds real, nameable structure (extractable back
into human mathematics, in the spirit of \cref{sec:hope}) or merely a
more efficient encoding of the same statistical residue is, at present,
an open and very testable question---and a natural apex to the
undergraduate program, where a student can point a modern learning tool
at a scaled-down family and see, concretely, what a machine can and
cannot pull out of manufactured randomness.

\begin{project}{A neural distinguisher at toy scale}{programming;
  a machine-learning course}{6--8 weeks (flagship)}
  Reproduce Gohr's experiment~\citep{Gohr2019} on a \SHA{}-family toy:
  train a small network to distinguish $r$-round outputs (or output
  differences) from random, and chart its advantage against $r$
  alongside the trail-toll curve of the miniature-trails project above
  and the statistical-battery frontier of \cref{sec:gap}. Then
  ask the hard question: \emph{interpret} the trained model---does its
  decision reduce to a differential-distribution feature the human
  theory already knows (\cref{sec:diffcalc}), or to something new? A
  negative result is a clean measurement; a positive one is the first
  thread of a natural property in a real hash family, and worth far
  more than a course grade.
\end{project}
